\section{Introduction}
\label{sec:introduction}

Modern foundation models have shifted the performance frontier of AI inference
from arithmetic throughput to \emph{data coordination}.
When model weights, activations, and key--value (KV) caches exceed accelerator
memory, inference latency is governed less by compute and more by how efficiently
execution coordinates data movement across a hierarchical memory system spanning
device memory, host memory, and secondary
storage~\cite{ma2026inference,gholami2024memorywall}.

Hardware analyses identify memory capacity, bandwidth, and interconnect
latency as important constraints on
inference~\cite{zhou2024survey}.
Separating capacity pressure from exposed transfer is therefore useful when
reporting and comparing inference systems.

Existing systems reduce memory pressure via CPU/NVMe
offloading~\cite{flexgen2023,aminabadi2022deepspeedinference},
GPU--CPU swapping~\cite{huang2020swapadvisor},
and KV-cache management~\cite{kwon2023pagedattention}.
Their outcomes depend on model size, batch size, hardware topology, and the
amount of traffic that remains on the critical path. This motivates a compact
description that exposes those conditions before comparing policies.

We answer this question with a regime-based view.
Denning's working-set model and thrashing theory~\cite{mattson1970evaluation}
identify capacity saturation as a binary event on a \emph{single-tier} system,
and the Roofline model~\cite{williams2009roofline} expresses performance as
$\min(F \cdot I,\, B_{\mathrm{peak}})$ for a fixed workload on a single tier.
Neither framework addresses a \emph{multi-tier} hierarchy in which computation,
locality, cross-tier transfer, and synchronisation co-exist as distinct,
adjustable latency terms.
Our accounting model decomposes end-to-end latency into four terms
($T_{\mathrm{comp}} + T_{\mathrm{mem}} + T_{\mathrm{swap}} + T_{\mathrm{sync}}$)
and uses two dimensionless ratios to label an operating point.
The three operational labels are governed by the
fast-memory residency ratio $R_C = C_{\mathrm{fast}}/W$ and the
overlap ratio $R_B = B_{\mathrm{slow}}T_{\mathrm{comp}}/D
= T_{\mathrm{comp}}/T_{\mathrm{transfer}}$.
The labels state whether majority residency is available and whether ideal
compute--transfer overlap can hide compulsory traffic; they do not by
themselves predict a policy's speed-up.

We make three contributions:
\begin{itemize}
  \item \textbf{Operational formulation.}
    We separate residency from compute--transfer overlap using two measurable,
    dimensionless ratios and state an explicit three-label classification rule.
  \item \textbf{Minimal analytical model.}
    A structural lower bound on the irreducible memory-access and transfer
    costs, with regime boundaries following directly from the two ratios:
    $\theta_B = 1.0$ derived as the overlap boundary
    ($R_B = T_{\mathrm{comp}}/T_{\mathrm{transfer}} = 1$), and $\theta_C = 0.50$
    as the majority-residency crossover.
  \item \textbf{Open proof of concept.}
    A compiled dependent-load CPU probe illustrates working-set sensitivity,
    and a separate layered-matrix harness exercises the measurement protocol.
    We distinguish these CPU observations and unexecuted CUDA/XLA paths from
    accelerator evidence.
\end{itemize}

Together, the formulation and CPU data provide a reproducible
starting point for testing when hierarchical-memory coordination is constrained
by residency or exposed transfer. They do not establish effect sizes or
strategy rankings on accelerators.

\FloatBarrier
